\documentclass[11pt]{article}
\usepackage[a4paper,margin=1in]{geometry}
\usepackage{amsmath,amssymb,amsthm,hyperref}
\newtheorem{theorem}{Theorem}
\newtheorem{lemma}{Lemma}
\newcommand{\Mat}{\operatorname{Mat}}
\newcommand{\GL}{\operatorname{GL}}
\newcommand{\rank}{\operatorname{rank}}
\title{Filling matrix spaces by transformed sets of arbitrary positive density exponent}
\author{Principia Math In Progress}
\date{22 September 2026}
\begin{document}
\maketitle
\begin{abstract}
We prove the matrix-set filling conjecture of Gill, Lifshitz, Pyber and
Szab\'o for every positive size exponent, every finite field and every
matrix dimension. The many-rank point-mass theorem proved in the companion
note, together with the source's additive-energy and growth lemmas,
gives an explicit number of summands depending only on the size exponent.
All small ambient spaces are included.
\end{abstract}

\section*{Task 1: the matrix-set filling theorem}
Write $A=\Mat_{n\times n}(\mathbb F_q)$ and $Q=|A|=q^{n^2}$, where
$q$ is a prime power and $n\ge1$. For subsets, sums denote Minkowski sums.
The following proves Conjecture 5.4 of \cite{source}.

\begin{theorem}\label{main}
For every $c>0$ there is an integer $\mu(c)$ such that for any
$X_1,\ldots,X_{\mu(c)}\subseteq A$ with $|X_i|\ge Q^c$, there exist
$a_i,b_i\in\GL(n,q)$ satisfying
\[
 a_1X_1b_1+\cdots+a_{\mu(c)}X_{\mu(c)}b_{\mu(c)}=A.
\]
For $0<c\le1$ one may take
\[
 \begin{gathered}
 D=2^{\max\{15,\lceil8/c\rceil\}},\qquad
 \ell=\lceil16/c\rceil,\qquad h=\ell+\lfloor4/c\rfloor,\\
 \nu=h2^{h+1},\qquad \mu(c)=2D\nu.
 \end{gathered}
\]
For $c>1$ the hypotheses are impossible, so $\mu(c)=1$ suffices.
\end{theorem}

The sets may be different and need not contain zero or be symmetric.
The source's established exponent $c=9/10$ is its Theorem 5.2,
equivalently Theorem 5.28; that case and the two lemmas used below are
prior results, not new claims of this note.

\section*{The inputs and their exact scope}
Let $T=\GL(n,q)\times\GL(n,q)$ act on the additive group $A$ by
$(g,h)x=gxh^{-1}$. This is an action by group automorphisms.
We use the following two results from the completely read primary source.

\begin{lemma}[\cite{source}, Corollary 5.13]\label{growth}
Let $B$ be a finite abelian group and $T\le\operatorname{Aut}(B)$.
For any nonempty subsets $Y_1,\ldots,Y_L$ there exist $s_i\in T$ such that
either
\[
 \left|\sum_{i=1}^L s_iY_i\right|\ge\frac L2\min_i|Y_i|,
\]
or that sum is a union of cosets of a $T$-invariant subgroup $N$
with $|N|\ge\min_i|Y_i|$.
\end{lemma}

\begin{lemma}[\cite{source}, Corollary 5.17]\label{energy}
Suppose $B$ is finite abelian, $T\le\operatorname{Aut}(B)$ is finite,
$\alpha,\beta\ge1$, and $V\subseteq B$ has $|V|\le\beta$.
Suppose for every $a\in B$ and $b_1,\ldots,b_\ell\notin V$,
\[
 \#\{(t_1,\ldots,t_\ell)\in T^\ell:
               a=t_1b_1+\cdots+t_\ell b_\ell\}\le |T|^\ell/\alpha.
\]
For $\delta>0$, set $h=\ell+\lfloor1/\delta\rfloor$ and $\nu=h2^{h+1}$.
If $Y_1,\ldots,Y_\nu$ have size at least $\alpha^\delta\beta$, then
there exist $s_i\in T$ such that their transformed sum either has size
at least $\alpha$, or contains a coset of a $T$-invariant subgroup $N$
with $|N|\ge\min_i|Y_i|$.
\end{lemma}

These inputs are proved in the source using Olson's growth theorem
and additive energy. In particular, they apply to unequal sets and to
cosets, and their parameters have no dependence on $|B|$.

The new input is the companion point-mass theorem \cite{rankbound}:
independent uniform matrices of ranks $r_1,\ldots,r_k$ have each point
probability at most $C/Q$, where $C=32768$, whenever $\sum_i r_i\ge2n$.
For clarity, its mechanism is rectangular factorization: the probability
that an $n\times r$ uniform factor is full rank is at least $4^{-r/n}$.
Choose a rank prefix with total $2n\le R<3n$. Conditioning both factors
in each block costs at most $16^3$, while the unconditioned product
of an $n\times R$ and an $R\times n$ uniform matrix has point mass
at most $8/Q$. Further independent summands do not increase the maximum.
The full proof, exact verification, and fixed publication version of
this dependency are linked in \cite{rankbound} and the accompanying audit.

\section*{Invariant subgroups and the exceptional ranks}
The only $T$-invariant additive subgroups of $A$ are $\{0\}$ and $A$.
Indeed, a nonzero invariant subgroup contains, by left-right equivalence,
$J_r=\operatorname{diag}(I_r,0)$ for some $r\ge1$.
When $n\ge2$, invariance and subtraction give
\[
 (I_n+E_{21})J_r-J_r=E_{21}.
\]
Its left-right orbit consists of every rank-one matrix. These include all
scalar multiples of matrix units and additively generate $A$.
When $n=1$, the orbit of any nonzero element already contains all
nonzero field elements. This also proves the assertion over nonprime fields
without assuming in advance that an additive subgroup is field-linear.

Set $\eta=c/8$ and let $V$ consist of matrices of rank less than $\eta n$.
The elementary rank-ball estimate in Fact 5.23(d) of \cite{source} gives
\[
 |V|\le4q^{2\eta n^2}=4Q^{c/4}.
\]
For completeness, if $s$ is the largest integer less than $\eta n$,
every rank-at-most-$s$ matrix has image in some $s$-dimensional subspace.
There are at most $4q^{s(n-s)}$ such subspaces and at most $q^{ns}$
maps into each, giving at most $4q^{s(2n-s)}\le4q^{2\eta n^2}$.
Here $s\ge0$ because $c>0$, and $\eta<1$ for the range under consideration.

For fixed $b_i\notin V$, independently uniform $t_i\in T$ send $b_i$
to independent uniform matrices of ranks at least $\eta n$.
Since $\ell=\lceil16/c\rceil$ gives $\ell\eta\ge2$, the point-mass
theorem implies
\[
 \#\{(t_i)\in T^\ell:\sum_i t_ib_i=a\}\le C|T|^\ell/Q.
                                                               \tag{1}
\]
Uniformity here is the elementary orbit-stabilizer fact; no independence
of matrix entries within a fixed rank class is assumed.

\section*{The two size regimes}
\begin{proof}[Proof of Theorem \ref{main}]
It remains to consider $0<c\le1$. Every $X_i$ has at least two elements,
since $Q^c>1$ and its cardinality is an integer.

First suppose $Q\ge D$. Apply Lemma \ref{energy} with
\[
 \alpha=Q/D\ge1,\qquad \beta=4Q^{c/4}\ge1,\qquad \delta=c/4.
\]
Its orbit-count hypothesis follows from (1), because $D\ge C$.
The size condition holds since $D^{c/4}\ge4$ and
\[
 \alpha^\delta\beta
 =4D^{-c/4}Q^{c/2}\le Q^{c/2}\le Q^c\le |X_i|.
\]
Partition the $\mu=2D\nu$ sets into $2D$ blocks of $\nu$ sets each.
Lemma \ref{energy} transforms each block into a sum $Y_j$ of size at
least $\alpha$. In its subgroup alternative, the subgroup has at least
two elements and therefore equals $A$, so that block already fills $A$.
In either case $|Y_j|\ge\alpha$, and also $|Y_j|\ge2$ because a
nonempty sumset contains a translate of each summand.

Now apply Lemma \ref{growth} to $Y_1,\ldots,Y_{2D}$.
Its first alternative gives size at least
$D\min_j|Y_j|\ge D\alpha=Q$. Its subgroup alternative also gives all
of $A$, since the subgroup has size at least two. Thus transformed
block sums fill $A$. Composing their outer transformations with the
inner transformations gives one element of $T$ for each original set.
Writing these as $a_iX_ib_i$ yields the required form, since an inverse
of an invertible right multiplier is again an allowed right multiplier.

If $Q<D$, apply Lemma \ref{growth} directly to the $\mu$ original sets.
The size alternative gives at least $\mu\ge Q$ elements, while the
subgroup alternative again forces the whole group because
$\min_i|X_i|\ge2$. This proves the result in every remaining dimension
and field, including all small ambient spaces.
\end{proof}

This is a separate set-filling assertion. Its proof explicitly uses the
point-mass theorem of \cite{rankbound}; that dependency is not claimed
as a second new point-mass result.

\section*{Proposed follow-up Tasks 2--5}
These are proposed tasks, not additional completed results.
\begin{enumerate}
\setcounter{enumi}{1}
\item Determine the optimal growth rate of the least admissible $\mu(c)$
as $c\downarrow0$, and compare it with the explicit bound above.
\item Give an algorithm for constructing the left and right multipliers,
with a complexity bound in terms of the input sets and field size.
\item Establish the corresponding filling theorem for rectangular matrix
spaces, with constants uniform in both dimensions and their ratio.
\item Characterize the subsets for which left multipliers alone suffice,
and identify all invariant-subspace obstructions to a uniform theorem.
\end{enumerate}
\begin{thebibliography}{9}
\bibitem{source} N. Gill, N. Lifshitz, L. Pyber and E. Szab\'o,
\emph{Initiating the proof of the Liebeck--Nikolov--Shalev conjecture},
\href{https://arxiv.org/abs/2408.07800v2}{arXiv:2408.07800v2}, 2024.
Conjecture 5.4, Corollaries 5.13 and 5.17, and Theorem 5.28.
\bibitem{rankbound} Principia Math In Progress,
\emph{A uniform point-mass bound for sums of prescribed-rank matrices},
22 September 2026, MathDB 364356, Theorem 1.
\href{https://github.com/antoshashakov/Principia-Math-In-Progress/blob/b01f62e992fc0a2038c41c3a728bbe8c928c1896/mathdb-open-problems/problems/364356/solution.pdf}{Fixed published proof version}.
\end{thebibliography}
\end{document}
